% Generated from validated Article Draft Result. Do not hand-edit; revise the structured draft instead.
\section{2022年に変わったのは、modelだけではなく設計空間だった}
\label{pkg:annual-synthesis}
\noindent\textbf{年初のinstruction followingとcompute allocation、春夏のmultimodal・media generationとmodel access、夏以降のopen distribution、秋のaction/retrievalとevaluation、年末のdialogue/controlをつなぐと、2022年は『巨大modelの競争』から『複数layerを組み合わせる生成AI systemの設計』へ重心が移った年として読める。ただし、2023年以降のagent化やplatform成熟をこの年へ逆投影してはいけない。}\autocite{src-16a1a1dc486e,src-1e362767e63c,src-26aadbc0393f,src-27297a80fac0,src-2832ee43ab8f,src-2c38d2f354fa,src-6ffa5c4e64ca,src-811e1008822e,src-86ce6205a4ed,src-87477ad8b476,src-929080260e16,src-bc1a61aec959,src-e58f1e9c8754,src-e8da8e8f02bc}

年初と年末を比較すると、問いの置き方が変わっている。InstructGPTはpretrained modelをどう人間の指示へ合わせるかを主要問題として可視化し、Chinchillaはscaleをtraining computeとdata配分まで含む設計問題として捉え直した。年末にはChatGPT research previewが、こうしたmodel behaviorの設計を利用者向け対話interfaceの上で見せる境界になった。modelそのものの能力だけでなく、post-trainingとinterfaceを含むsystem behaviorが競争軸になったのである。\autocite{src-6ffa5c4e64ca,src-87477ad8b476,src-bc1a61aec959,src-e58f1e9c8754,src-e8da8e8f02bc}


別の軸では、DALL-E 2とStable Diffusionが生成メディアを研究成果だけでなく異なるdistribution surfaceへ展開し、ReActはreasoningを外部actionへ接続する方向を示した。HELMのようなevaluation frameworkが必要になった背景には、modelがtext completionだけではなくmedia、code、action、dialogueなど異質なsystemへ広がり、一つのbenchmarkでは比較しにくくなったという構造変化がある。distribution、interface、evaluationはmodel本体の外側にあるが、生成AIの実用性を決めるlayerとして無視できなくなった。\autocite{src-16a1a1dc486e,src-1e362767e63c,src-26aadbc0393f,src-27297a80fac0,src-2832ee43ab8f,src-2c38d2f354fa,src-811e1008822e,src-86ce6205a4ed,src-929080260e16}


変わったのは、設計のレバーが増え、それぞれを独立して議論できるようになったことである。pretraining scale、post-training、prompt-time reasoning、compute/data allocation、model access、runtime efficiency、media distribution、action/retrieval、evaluation、controlが別々の層として見えるようになった。一方で、performance claimの独立再現、openという語の定義、長期的な安全性、reward objectiveの限界、異なるbenchmark間の比較可能性といった制約は解消されていない。layerが増えたことは、制約が消えたことを意味しない。\autocite{src-1e362767e63c,src-26aadbc0393f,src-2832ee43ab8f,src-6ffa5c4e64ca,src-87477ad8b476,src-e58f1e9c8754,src-e8da8e8f02bc}


2022年を評価するとき、2023年以降の結果を知っていることは強い誘惑になる。ChatGPTの後年のplatform化、agent/tool useの一般化、open-weight ecosystemの拡大、生成videoの成熟などは、2022年時点ではまだ同じ形で成立していない。本版では、後年に重要になったtrajectoryを手掛かりにはしても、その成熟状態を2022年のartifactへ逆投影しない。2022年に確認できるのは、後のsystem化につながるlayerが年内に分離・接続され始めたことまでである。\autocite{src-26aadbc0393f,src-2832ee43ab8f,src-811e1008822e,src-bc1a61aec959}


したがって2022年の核心は、生成AIが一つの巨大modelを中心に説明できる領域ではなくなったことである。能力を作るpretraining、振る舞いを整えるsteering、計算とaccessを支えるsystems、mediaを配布するecosystem、外部世界へ接続するinterface、強いsystemを測るevaluation、利用者向け挙動を制御するcontrolが相互依存する構造へ移った。ChatGPTはその変化の唯一の原因ではなく、複数layerが対話UI上で一度に見えるようになった年末の象徴として位置付けるのが、本版の年次結論である。\autocite{src-16a1a1dc486e,src-1e362767e63c,src-26aadbc0393f,src-27297a80fac0,src-2832ee43ab8f,src-2c38d2f354fa,src-6ffa5c4e64ca,src-811e1008822e,src-86ce6205a4ed,src-87477ad8b476,src-929080260e16,src-bc1a61aec959,src-e58f1e9c8754,src-e8da8e8f02bc}


\begin{claimboundary}
この総括は、各一次資料で検証した範囲を年次trajectoryへ圧縮した編集上の推論である。vendor・project・authorの性能主張を独立再現済みの事実へ変換せず、2022年に観測できる構造変化と後年から見た解釈を区別する。\autocite{src-1e362767e63c,src-87477ad8b476,src-e58f1e9c8754}
\end{claimboundary}
